\documentclass[14pt aspectratio=169]{beamer}


%%% Theme (color order, page numbering, etc)

\usetheme{Madrid}  % Copenhagen doesnt have page numbers; better dont change

\useoutertheme{split} 
% infolines - Автор (университет) | название | дата | число слайдов
% split - Автор | название | число слайдов + headline
% miniframes - infolines + headline


%%% Circle theme without shadows (uncomment latter, comment former)

\useinnertheme{circles}

%\setbeamertemplate{blocks}[rounded][shadow=false]
%\setbeamertemplate{items}[circle]
%\setbeamertemplate{sections/subsections in toc}[circle]
%\setbeamertemplate{title page}[default][colsep=-4bp,rounded=true,shadow=false]
%\setbeamertemplate{part page}[default][colsep=-4bp,rounded=true,shadow=false]
%\setbeamertemplate{section page}[default][colsep=-4bp,rounded=true,shadow=false]
%\setbeamertemplate{subsection page}[default][colsep=-4bp,rounded=true,shadow=false]


%%% Kill headline and navigation symbols

\setbeamertemplate{headline}{} % kill headline with section names
\setbeamertemplate{navigation symbols}{}  % kill navigation symbols


%%% colors

\definecolor{AUred}{rgb}{0.609 0 0.297}
\definecolor{AUblue}{rgb}{0.398 0 0.609}
%\definecolor{AUblue}{rgb}{0.199 0 0.398} darker blue

\setbeamercolor{palette primary}{bg=AUred,fg=white}
\setbeamercolor{palette secondary}{bg=AUred,fg=white}
\setbeamercolor{palette tertiary}{bg=AUred,fg=white}
\setbeamercolor{palette quaternary}{bg=AUblue,fg=white}
\setbeamercolor{structure}{fg=AUblue} % itemize, enumerate, etc
% But!!:
\setbeamercolor{caption name}{fg=black}
\setbeamercolor{section in toc}{fg=AUblue} % TOC sections

% Override palette coloring with secondary

\setbeamercolor{subsection in head/foot}{bg=AUred,fg=white}

%%% Misc

\setbeamercovered{transparent}  % make covered items opaque
\usepackage{appendixnumberbeamer}  % to hide slides in \appendix
\setbeamertemplate{caption}[numbered]

%%% Non-beamer packages

\usepackage[english, russian]{babel}
\usepackage[utf8]{inputenc}  % fonts for xelatex
%\setmainfont{arial}

\usefonttheme{serif}  % beamer way to set font

\usepackage{amsmath, amsfonts, amssymb, mathtools}  % mathfonts
\usepackage[f]{esvect}  % decent vector sign and double-dot command (\ddot{})
\usepackage{pifont}  % special symbols (bullets)

\usepackage{indentfirst}

\usepackage{array}
\usepackage{graphicx}
\usepackage{float}
\usepackage{wrapfig}
\usepackage{tikz}

\usepackage{multirow}
\newcolumntype{O}{>{$}c<{$}}  % mathmode in centered columns

\usepackage[font=scriptsize,labelfont=bf]{caption}
\usepackage{subcaption}
\setlength{\parindent}{20pt}

\theoremstyle{plain}
\newtheorem{thm}{Теорема}
\newtheorem{prf}{Доказательство}
\newtheorem{clr}{Следствие}
\newtheorem{lm}{Лемма}[]
\newtheorem*{lm*}{Лемма}
\newtheorem*{dfn}{Определение}
\newtheorem*{rem}{Замечание}
\newtheorem{ex}{Пример}
\newtheorem*{prop}{Предложение}
\newtheorem*{thm*}{Теорема}
\newtheorem*{ex*}{Пример}
\newtheorem*{clr*}{Следствие}
\newtheorem*{ill}{Пример}

%\newcommand{\sh}{\mathrm{sh} \,}
%\newcommand{\ch}{\mathrm{ch} \,}
%\newcommand{\th}{\mathrm{th} \,}
%\newcommand{\cth}{\mathrm{cth} \,}

%\newcommand{\arcsh}{\mathrm{arcsh} \,}
\newcommand{\arcch}{\mathrm{arcch} \,}
%\newcommand{\arcth}{\mathrm{arcth} \,}
%\newcommand{\arccth}{\mathrm{arccth} \,}


\usepackage{hyperref}
\hypersetup{
	colorlinks=false  % dont color links, cause name in the footline is link
}


\title{Задача №9 Хаотичные шнурки}
\author{Докладчик: Кононов Александр}
\institute{{\normalsize Команда: Не на паре, Алферовский университет}}
%\institute{Алферовский университет}
\date{}  % \today

\begin{document}

\begin{frame}
	\tikz [remember picture,overlay]
		\node at
		([xshift=-2cm, yshift=2cm]current page.south east)
		%or: (current page.center)
		{\includegraphics[width=.3\textwidth,height=.2\textheight]{img/au-logo-shadow.pdf}};
	\titlepage
\end{frame}


\section{Условие задачи}

\begin{frame}{Условие задачи}
		\begin{block}{Хаотичные шнурки}
			Если достаточно большое количество шнурков вращать в центрифуге с водой, подобно тому, как это происходит в стиральной машине, то шнурки запутаются. Исследуйте количество узлов на одном шнурке, а также узлы между шнурками. Рассмотрите самую длинную веревку, получающуюся в результате каждого такого опыта. Исследуйте распределение ее длины для серии таких опытов.
			%\textit{Подсказка: } Обычно это явление наблюдается при использовании жесткой воды для чая.
		\end{block}
	\begin{figure}
		\includegraphics[width=0.3\linewidth]{img/hnurki_logo.png}
	\end{figure}

\end{frame}

\begin{frame}{Что такое узел?}

	\begin{columns}

		\begin{column}{0.5\linewidth}
			\begin{block}{Определение узла}
				Назовем узлом замкнутую 2d нить устойчивую к 3 видам преобразований. То есть преобразованные нити являются топологически эквивалентными.
			\end{block}

			\begin{block}{Топология}
				Топология узла описывается полиномами Джинса.
			\end{block}

%			\begin{alertblock}{}
%				$\Rightarrow$ ячейки Рэлея-Бенара
%			\end{alertblock}
		\end{column}

		\begin{column}{0.5\linewidth}
			\includegraphics[height=1\linewidth]{img/hnurki_peracii.png}
		\end{column}

	\end{columns}

\end{frame}

\begin{frame}{Экспериментальная установка}

	\begin{block}{}
		Экспериментальная установка состоит из цилиндра длиной $l_0 = 13$см и двигателя с частотой вращения $\omega = 1 $ Гц
	\end{block}

	\begin{columns}

		\begin{column}{0.4\linewidth}
			\begin{figure}[H]
				\includegraphics[width=1\linewidth]{img/experiment.jpg}
				%\caption{Капля левитирует в восходящем потоке пара}
			\end{figure}

		\end{column}

		\begin{column}{0.6\linewidth}

			\begin{figure}[H]
				\includegraphics[width=1\linewidth]{img/exp_4.png}
				%\caption{Капля левитирует в восходящем потоке пара}
			\end{figure}
			Аппроксимация точек зависимостью:
			$P = \frac{P_0}{1+(l/l_0)^b}$
			где $P_0 = 0.3$, $b = -3$

		\end{column}

	\end{columns}

%	\begin{columns}
%
%		\begin{column}{0.5\linewidth}
%			\begin{figure}[H]
%				\includegraphics[width=\linewidth]{img/setting_exp.pdf}
%				\caption{Схема установки}
%			\end{figure}
%
%		\end{column}
%
%		\begin{column}{0.5\linewidth}
%
%			\begin{figure}[H]
%				\includegraphics[width=0.8\linewidth]{img/new_prem.pdf}
%				\caption{Фото узора из капель}
%			\end{figure}
%			\vspace{-0.5cm}
%			\begin{figure}[H]
%			\includegraphics[width=0.8\linewidth]{img/micro_ph.pdf}
%			\caption{Фото через дальнофокусный микроскоп}
%			\end{figure}
%
%		\end{column}
%
%	\end{columns}

\end{frame}

\begin{frame}{Теоретическое моделирование}

	\begin{block}{}
		В качестве модели взят марковский процесс случайного блуждания конца шнурка.
		Конец шнурка с равной вероятностью мог за один такт опуститься вверх или вниз над или под одной из частей шнурка.
		После случайного блуждания образуется узел.
	\end{block}

	\begin{columns}

		\begin{column}{0.4\linewidth}
			\begin{figure}[H]
				\includegraphics[width=1\linewidth]{img/hnurki_zaput.png}
				%\caption{Капля левитирует в восходящем потоке пара}
			\end{figure}

		\end{column}

%		\begin{column}{0.5\linewidth}
%
%			\begin{figure}[H]
%				\includegraphics[width=0.5\linewidth]{img/probobility_31.png}
%				%\caption{Фото узора из капель}
%			\end{figure}
%		\end{column}
%
%		\begin{column}{0.3\linewidth}
%			\begin{figure}[H]
%				\includegraphics[width=0.5\linewidth]{img/probobility_32.png}
%				%\caption{Фото через дальнофокусный микроскоп}
%			\end{figure}
%
%		\end{column}

	\end{columns}

\end{frame}

\begin{frame}{Результаты моделирования}
		\begin{block}{}
			Графики зависимости вероятности образования узла на одном шнурке от его длины и времени блуждания.
		\end{block}

%			\begin{figure}[H]
%				\includegraphics[width=0.7\linewidth]{img/model_1.png}
%				\caption{Необходимое число пересечений: 3}
%			\end{figure}


	\begin{columns}

		\begin{column}{0.5\linewidth}
			\begin{figure}[H]
				\includegraphics[width=1\linewidth]{img/time.png}
				%\caption{Необходимое число пересечений: 3}
			\end{figure}

		\end{column}

		\begin{column}{0.5\linewidth}
			\begin{figure}[H]
				\includegraphics[width=1\linewidth]{img/model_1.png}
				Необходимое число пересечений: 3
			\end{figure}
			Аппроксимация точек:
			$P = \frac{P_0}{1+(l/l_0)^b}$
			где $P_0 = 0.8$, $b = -3$
		\end{column}

%		\begin{column}{0.5\linewidth}
%			\begin{figure}[H]
%				\includegraphics[width=1\linewidth]{img/model_7.png}
%				\caption{Необходимое число пересечений: 7}
%			\end{figure}
%		\end{column}

%		\begin{column}{0.5\linewidth}
%
%			\begin{figure}[H]
%				\includegraphics[width=0.5\linewidth]{img/probobility_31.png}
%				%\caption{Фото узора из капель}
%			\end{figure}
%		\end{column}
%
%		\begin{column}{0.3\linewidth}
%			\begin{figure}[H]
%				\includegraphics[width=0.5\linewidth]{img/probobility_32.png}
%				%\caption{Фото через дальнофокусный микроскоп}
%			\end{figure}
%
%		\end{column}

	\end{columns}

\end{frame}

\begin{frame}{Результаты моделирования}
		\begin{block}{}
			Графики зависимости вероятности образования узла на одном шнурке от его длины с разными условиями на число пересечений в модели для образования узла.
		\end{block}

%			\begin{figure}[H]
%				\includegraphics[width=0.7\linewidth]{img/model_1.png}
%				\caption{Необходимое число пересечений: 3}
%			\end{figure}


	\begin{columns}

		\begin{column}{0.5\linewidth}
			\begin{figure}[H]
				\includegraphics[width=1\linewidth]{img/model_55.png}
				Необходимое число пересечений: 5
			\end{figure}
			Аппроксимация точек:
			$P = \frac{P_0}{1+(l/l_0)^b}$
			где $P_0 = 0.57$, $b = -3$
		\end{column}

		\begin{column}{0.5\linewidth}
			\begin{figure}[H]
				\includegraphics[width=1\linewidth]{img/model_7.png}
				Необходимое число пересечений: 7
			\end{figure}
			Аппроксимация точек:
			$P = \frac{P_0}{1+(l/l_0)^b}$
			где $P_0 = 0.33$, $b = -3$
		\end{column}
	\end{columns}

\end{frame}

\begin{frame}{Сравнение со статьей}

	\begin{block}{}
		В статье "Spontaneous knotting of an agitated string". Dorian M. Raymer and Douglas E. Smith представлены следующие результаты
	\end{block}

	\begin{columns}

		\begin{column}{0.4\linewidth}
			\begin{figure}[H]
				\includegraphics[width=1\linewidth]{img/knot_prob.png}
				%\caption{Капля левитирует в восходящем потоке пара}
			\end{figure}
			$P = \frac{P_0}{1+(l/l_0)^b}$
			\\где $P_0 = 0.55$, $b = -2.9$
		\end{column}

		\begin{column}{0.6\linewidth}
			\begin{figure}[H]
				\includegraphics[width=1\linewidth]{img/both_7.png}
				%\caption{Капля левитирует в восходящем потоке пара}
			\end{figure}
			$\, \, \, \, P = \frac{P_0}{1+(l/l_0)^b}$
			где $P_0 = 0.33$, $b = -3$
		\end{column}


	\end{columns}

\end{frame}

%\begin{frame}{Левитация капли}
%	\begin{figure}[H]
%		\includegraphics[width=0.6\linewidth]{img/droplet.pdf}
%		\caption{Капля левитирует в восходящем потоке пара}
%	\end{figure}
%\end{frame}


%\begin{frame}{Левитация капли}
%
%	$$
%	6 \pi u \mu R \cdot \beta \left( \dfrac{z}{R} \right)
%	= \dfrac{4 \pi}{3} R^3 \rho_w g
%	$$
%
%	где $\beta$ - поправка к закону Стокса ($\xi = \text{arc ch} \dfrac{z}{R}$):
%
%	\[
%	\beta \left( \dfrac{z}{R} \right)
%	= \dfrac{4}{3}
%	 \sh \xi
%	 \sum_{n=1}^{\infty}
%	 \dfrac{n(n+1)}{(2n-1)(2n+3)}
%	 \left[
%	 \dfrac{4 \ch^2 (n+\dfrac{1}{2}) \xi + (2n+1)^2 \sh^2 \xi}{2 \sh (2n+1)\xi - (2n+1) \sh 2\xi} -1
%	 \right]
%	\]
%
%	\begin{minipage}{0.4\linewidth}
%		\begin{tabular}{|O|O|}
%			\hline
%			z/R & \beta \\
%			\hline
%			1 & \infty \\
%			\hline
%			1.1 & 4.626 \\
%			\hline
%			1.2 & 3.104 \\
%			\hline
%			1.3 & 2.541 \\
%			\hline
%			1.4 & 2.235 \\
%			\hline
%		\end{tabular}
%		\begin{tabular}{|O|O|}
%			\hline
%			z/R & \beta \\
%			\hline
%			1.5 & 2.039 \\
%			\hline
%			2 & 1.597 \\
%			\hline
%			5 & 1.175 \\
%			\hline
%			10 & 1.081 \\
%			\hline
%			\infty &  1 \\
%			\hline
%		\end{tabular}
%	\end{minipage}
%	\begin{minipage}{0.4\linewidth}
%		{\footnotesize \label{happel} [\ref{happel}] J. Happel, H. Brenner, Low Reynolds number	hydrodynamics with special applications	to particulate media,  - Prentice-Hall, 1965}
%	\end{minipage}
%
%\end{frame}

%\begin{frame}{Левитация капли}
%	\begin{columns}
%		\begin{column}{0.55\linewidth}
%			\begin{figure}[H]
%				\includegraphics[width=\linewidth]{lev/hr_beta.pdf}
%				\caption{Зависимость высоты $h$ левитации капли от ее радиуса, $u$ - скорость вертикального потока (м/с)}
%			\end{figure}
%		\end{column}
%		\begin{column}{0.45\linewidth}
%			\[
%			\dfrac{ \beta \left( \dfrac{z}{R} \right) }{R^2}
%			=
%			\dfrac{2 \rho_w g}{9 u \mu}
%			\rightarrow \boxed{z = z (R)}
%			\]
%			\[
%			\mu \approx 1.16 \cdot 10^{-5} ~ \dfrac{kg}{m~s}, ~ \rho_w \approx 10^3 \dfrac{kg}{m^3}
%			\]
%		\end{column}
%	\end{columns}
%
%	$u$ и $z/R$ определим на эксперименте и проверим модель
%\end{frame}

%\begin{frame}{Структура пленки}
%
%	\begin{alertblock}{}
%		Пленка - средство диагностики процессов, происходящих в тонком слое жидкости, близком к ее поверхности
%	\end{alertblock}
%
%	\begin{figure}[H]
%			\includegraphics[width=0.8\linewidth]{img/ienna_cells.jpg}
%			\caption{Side-by-side comparison between the emitted infrared mosaic
%				pattern (a) and the visible distribution of floating droplets (b). Several of
%				the cells in both images are highlighted in yellow to facilitate the
%				similarity, {\label{ienna} [\ref{ienna}] - Spatially resolved evaporative patterns from water, Ienna F. et al}}
%	\end{figure}
%
%\end{frame}

%\begin{frame}{Структура пленки}
%
%	\begin{columns}
%		\begin{column}{0.6\linewidth}
%			\begin{figure}[H]
%				\includegraphics[width=\linewidth]{img/ienna_infrared.jpg}
%				\caption{Infrared image of the surface of water, viewed from directly
%					above. Corresponding temperatures are shown on the right. [\ref{ienna}]}
%			\end{figure}
%		\end{column}
%
%		\begin{column}{0.2\linewidth}
%			Число Релея:
%			\[
%			Ra=\dfrac{g \beta \Delta T L^3}{\nu \chi}
%			\]
%		\end{column}
%	\end{columns}
%
%
%\end{frame}

%\begin{frame}{Что происходит в кружке?}
%
%	\begin{columns}
%
%		\begin{column}{0.5\linewidth}
%			\includegraphics[height=1\linewidth]{img/br_set.pdf}
%		\end{column}
%
%		\begin{column}{0.5\linewidth}
%			\begin{block}{Толстый слой}
%				Теплопроводные непроницаемые стенки и дно; граница с тонким слоем непроницаемая, темепратура меняется непрерывно
%			\end{block}
%
%			\begin{block}{Тонкий слой}
%				Бесконечная поверхность жидкости, верхняя границе открыта, считаем что воздух тепло не проводит
%			\end{block}
%
%			\begin{alertblock}{}
%				$\Rightarrow$ ячейки Рэлея-Бенара
%			\end{alertblock}
%		\end{column}
%
%	\end{columns}
%
%\end{frame}

%\begin{frame}{Образование страйпа}
%
%	\begin{columns}
%
%		\begin{column}{0.6\linewidth}
%			\centering
%			\begin{figure}
%				\includegraphics[width=\linewidth]{str/denser.jpg}
%				\caption{На фотографии ([\ref{ienna}]) видно, что плотность капель пара выше вблизи границ ячеек - там и формируются страйпы}
%			\end{figure}
%		\end{column}
%
%		\begin{column}{0.4\linewidth}
%
%			\begin{block}{}
%				\begin{itemize}
%					\item Высокая плотность
%					\item Малая высота
%					\item Нисходящие потоки
%				\end{itemize}
%			\end{block}
%
%		\end{column}
%
%	\end{columns}
%
%	Капли сразу над поверхностью в силу вязкости воздуха и пара до некоторой степени повторяют движение потоков жидкости на поверхности жидкости. Движение этих потоков - движение от центра ячейки к ее краю, поэтому капли скапливаются на периферии, и при достижении достаточно большой концентрации возникает капиллярная волна
%
%\end{frame}

%\begin{frame}{Капилярная волна}
%	\begin{columns}
%		\begin{column}{0.6\linewidth}
%			\begin{figure}[H]
%				\includegraphics[width=\linewidth]{img/cap_draft.pdf}
%				\caption{Механизм капилярной волны}
%			\end{figure}
%		\end{column}
%		\begin{column}{0.4\linewidth}
%			$1\rightarrow2$: одна из капель коалесцирует с поверхностью воды
%
%			$2\rightarrow3$: капля растекается
%
%			$3\rightarrow4$: из-за искривленния поверхности, коалесцирует вторая капля
%		\end{column}
%	\end{columns}
%
%
%\end{frame}

%\begin{frame}{Жесткость}
%	\begin{columns}
%		\begin{column}{0.7\linewidth}
%			\begin{block}{}
%				Мягкая вода: $(0 \div 7 ~ \text{мг-экв/л}) \leftrightarrow (0 \div 0.14) \%$
%
%				Жесткая вода:  $(10 \div 20 ~ \text{мг-экв/л}) \leftrightarrow (0.2 \div 0.4) \%$
%			\end{block}
%		\end{column}
%		\begin{column}{0.2\linewidth}
%			\begin{block}{}
%				$t=d^2/D$
%			\end{block}
%		\end{column}
%	\end{columns}
%
%	\begin{figure}[H]
%		\includegraphics[width=0.5\linewidth]{sol/fedorets_sol.jpg}
%		\caption{Equilibrium water droplets in stabilized droplet clusters: domain of self-stabilized droplet clusters; [2] - Fedorets et al}
%	\end{figure}
%	\vspace{-0.5cm}
%	{\small Что подтвержает утвержджение в подсказке: образование структуры легче наблюдать в жесткой воде}
%\end{frame}

\begin{frame}{Вывод}
	\begin{itemize}
		%\item Качественно объяснена причина возникновения узлов
		%\item Рассмотрен механизм образования узлов на шнурках
		\item Экспериментально измерена вероятность образования узла на одном шнурке $P=0.3$
		\item Построена модель и полученная вероятность образования узла на одном шнурке $P=0.3$
		%\item Полученные результаты не сходятся между собой и со статьей
		\item В качестве доработки задачи можно рассмотреть распределения узлов между несколькими веревками, как модель случайного блуждания концов рядом с образовавшимися узлами на соседних веревках, распределенными равномерно по длине шнурка
	\end{itemize}
\end{frame}

\begin{frame}{Спасибо за внимание}
	Благодарим за внимание!
\end{frame}

%\appendix
%
%\begin{frame}{Скрытые слайды}
%	Дальше будут скрытые слайды...
%\end{frame}

%\begin{frame}{Поправка}
%
%	\begin{figure}[H]
%		\includegraphics[width=0.8\linewidth]{lev/beta.png}
%		\caption{Поправка к закону Стокса}
%	\end{figure}
%
%
%\end{frame}
%
\end{document}
%
%\begin{frame}{Испарение}
%
%	\begin{columns}
%		\begin{column}{0.5\linewidth}
%			\begin{figure}[H]
%				\includegraphics[width=\linewidth]{lev/evap.pdf}
%				\caption{Частица вылетает из воды в пар. Температуры на границе непрерывна}
%			\end{figure}
%		\end{column}
%
%		\begin{column}{0.5\linewidth}
%
%			\begin{block}{Начальная скорость}
%				$$
%				u_0 = \dfrac{\pi}{8}  v_T = \dfrac{\pi}{8} \sqrt{5RT/\mu}
%				$$
%				$v_T = \sqrt{5RT/\mu}$ - скорость теплового движения частиц воды
%			\end{block}
%
%		\end{column}
%	\end{columns}
%
%\end{frame}
%
%\begin{frame}{Восходящий поток}
%
%	\begin{columns}
%
%		\begin{column}{0.5\linewidth}
%			\begin{figure}[H]
%				\includegraphics[width=0.8\linewidth]{lev/colis.pdf}
%				\caption{Столкновение двух частиц. Частица сверху неподвижна, частица снизу имеет вертикальную скорость $u$}
%			\end{figure}
%		\end{column}
%
%		\begin{column}{0.5\linewidth}
%			\begin{block}{За одно столкновение}
%				$$
%				u' = u \cdot \dfrac{1}{\pi} \int_{-\pi/2}^{\pi/2} \sin^2 \theta d \theta = \dfrac{u}{2}
%				$$
%			\end{block}
%
%			\begin{block}{}
%				$$
%				u(z) = u_0 \cdot 2^{-z/l}, ~~ l = \lambda \dfrac{2}{\pi}
%				$$
%				$\lambda $ - длина свободного пробега частицы в паре
%
%			\end{block}
%
%
%		\end{column}
%	\end{columns}
%
%\end{frame}

%\begin{frame}{Скорость восходящего потока}
%
%	\begin{columns}
%		\begin{column}{0.7\linewidth}
%			\begin{figure}[H]
%				\includegraphics[width=\linewidth]{lev/rasch.png}
%				\caption{Скорость потока в зависимости от высоты над поверхностью воды $u(z)$}
%			\end{figure}
%		\end{column}
%
%		\begin{column}{0.3\linewidth}
%			\begin{block}{}
%				На масштабе порядка размера капли ($7$ мкм) скорость потока меняется сильно
%			\end{block}
%		\end{column}
%	\end{columns}
%
%\end{frame}
